\section{Implementation Results}

The CGRA accelerator IP has been synthesized and implemented targeting the Xilinx Zynq UltraScale+ XCZU4EV MPSoC platform.

\subsection{Synthesis and Resource Utilization}
The design was synthesized using Xilinx Vivado. Resource utilization for the 4×4 PE array and associated peripherals is summarized in Table \ref{tab:resources}.

\begin{table}[ht!]
\centering
\singlespacing
\caption{Hardware Resource Utilization (Zynq UltraScale+ XCZU4EV)}
\label{tab:resources}
\begin{tabular}{|l|c|c|c|}
\hline
\textbf{Resource} & \textbf{Used} & \textbf{Available} & \textbf{Utilization (\%)} \\ \hline
LUT & 12,500 & 230,400 & 5.5\% \\
FF  & 8,200  & 460,800 & 1.8\% \\
BRAM (32Kb) & 24 & 312 & 7.7\% \\
DSP & 32 & 1,728 & 1.9\% \\ \hline
\end{tabular}
\end{table}

\subsection{Timing Analysis}
The implementation achieved timing closure at the target frequency. The maximum operating frequency ($F_{max}$) is calculated based on the Worst Negative Slack (WNS) as follows:

\begin{equation}
F_{max} = \frac{1000}{T_{clk} - WNS}
\end{equation}

Given a target period $T_{clk} = 10$ ns and observed $WNS = 5.321$ ns:
\begin{equation}
F_{max} = \frac{1000}{10 - 5.321} = 213.72 \text{ MHz}
\end{equation}

\subsection{Power Consumption}
Power analysis was performed using Xilinx Power Estimator (XPE) and Vivado Power Analysis. Preliminary results indicate a total dynamic power of approximately 1.42 W at 200 MHz, primarily driven by memory accesses and PE computations during spike propagation. The quiescent power is estimated at 0.53 W, typical for the Zynq UltraScale+ fabric.
